\documentclass[../general/general]{subfiles}

\begin{document}

\chapter{Лекция №7}
\noindent\textit{А.\,М.~Белемук \hfill 27 марта 2021}
\vspace{1em}


% ===== СЕКЦИЯ 1: Свойства бозе-интегралов =====
\section{Свойства бозе-интегралов}

Итак, лекция седьмая, статистическая физика. На прошлой лекции мы вывели основные формулы для идеального бозе-газа: ввели распределение Бозе-Эйнштейна и записали уравнение на среднее число частиц через так называемые бозе-интегралы. Сегодня начнём с исследования свойств этих интегралов, а затем перейдём к важнейшему явлению --- бозе-эйнштейновской конденсации.

\subsection{Определение бозе-интегралов}

Напомним введённые на прошлой лекции обозначения. Безразмерный параметр $y = -\mu/T > 0$ (поскольку $\mu < 0$ для бозе-газа). Бозе-интегралы определяются как:
\begin{equation}
    I_\alpha(y) = \int_0^\infty \frac{x^\alpha}{e^{x+y} - 1}\, dx, \quad y > 0, \quad x = \frac{\epsilon}{T}, \quad y = -\frac{\mu}{T}
    \label{eq:bose_integral_def}
\end{equation}

Тогда уравнение на число частиц записывается в виде:
\begin{equation}
    N = A\,T^{3/2}\, I_{1/2}(y)
    \label{eq:N_bose}
\end{equation}
где $A$ --- коэффициент в плотности состояний $\nu(\epsilon) = A\sqrt{\epsilon}$. Нам будут нужны прежде всего $I_{1/2}(y)$ и $I_{3/2}(y)$.

\subsection{Свойство 1: монотонное убывание}

Первое свойство: $I_\alpha(y)$ монотонно убывает по $y$ при $y > 0$.

Это видно непосредственно из определения: при увеличении $y$ аргумент экспоненты $e^{x+y}$ возрастает, знаменатель увеличивается, а значит подынтегральная функция уменьшается в каждой точке. Следовательно, и весь интеграл убывает.

\subsection{Свойство 2: значения при $y = 0$}

Второе свойство: при $y = 0$ бозе-интегралы принимают конечные значения, выражающиеся через гамма-функцию и дзета-функцию Римана:
\begin{equation}
    I_\alpha(0) = \int_0^\infty \frac{x^\alpha}{e^x - 1}\, dx = \Gamma(\alpha+1)\,\zeta(\alpha+1)
\end{equation}

Это стандартный интеграл, который вычисляется разложением в ряд. Действительно, при $y = 0$:
\begin{equation}
    \frac{1}{e^x - 1} = \frac{e^{-x}}{1 - e^{-x}} = \sum_{k=1}^\infty e^{-kx}
\end{equation}
Подставляем:
\begin{equation}
    I_\alpha(0) = \sum_{k=1}^\infty \int_0^\infty x^\alpha e^{-kx}\, dx = \sum_{k=1}^\infty \frac{\Gamma(\alpha+1)}{k^{\alpha+1}} = \Gamma(\alpha+1)\,\zeta(\alpha+1)
\end{equation}

Для нужных нам значений $\alpha$:
\begin{align}
    I_{1/2}(0) &= \Gamma\!\left(\tfrac{3}{2}\right)\zeta\!\left(\tfrac{3}{2}\right) = \frac{\sqrt{\pi}}{2}\cdot\zeta\!\left(\tfrac{3}{2}\right) \approx \frac{\sqrt{\pi}}{2}\cdot 2{,}612 \approx 2{,}315 \\[4pt]
    I_{3/2}(0) &= \Gamma\!\left(\tfrac{5}{2}\right)\zeta\!\left(\tfrac{5}{2}\right) = \frac{3\sqrt{\pi}}{4}\cdot\zeta\!\left(\tfrac{5}{2}\right) \approx \frac{3\sqrt{\pi}}{4}\cdot 1{,}341 \approx 1{,}783
\end{align}

Запомним числовые значения: $\zeta(3/2) \approx 2{,}612$ и $\zeta(5/2) \approx 1{,}341$.

\subsection{Свойство 3: классический предел}

При $y \gg 1$ (высокие температуры, $\mu \to -\infty$) экспонента $e^{x+y} \gg 1$ для любого $x \geq 0$, поэтому единицей в знаменателе можно пренебречь:
\begin{equation}
    I_\alpha(y) \approx \int_0^\infty x^\alpha e^{-(x+y)}\, dx = e^{-y}\int_0^\infty x^\alpha e^{-x}\, dx = \Gamma(\alpha+1)\,e^{-y}
\end{equation}

Это классический предел. Убывание идёт как $e^{-y} = e^{\mu/T}$, что соответствует больцмановскому фактору Гиббса.

\subsection{Свойство 4: поведение производных при $y \to 0$}

Четвёртое, более тонкое свойство касается аналитичности функций $I_\alpha(y)$ при $y \to 0^+$.

Оказывается, $I_{1/2}(y)$ \emph{не аналитична} в нуле: её производная по $y$ обращается в бесконечность при $y \to 0^+$. Это можно увидеть из поведения $\frac{d}{dy}I_{1/2}(y) = -I_{-1/2}(y)$, интеграл $I_{-1/2}(y)$ логарифмически расходится при $y \to 0$.

В то же время $I_{3/2}(y)$ аналитична: $\frac{d}{dy}I_{3/2}(y) = -I_{1/2}(y)$, и $I_{1/2}(0)$ конечно. Следовательно, $I_{3/2}(y)$ имеет конечную производную при $y \to 0$.

Это различие имеет важные физические последствия: теплоёмкость бозе-газа будет иметь излом при температуре конденсации, а не настоящий фазовый переход второго рода (хотя для трёхмерного газа производная теплоёмкости претерпевает скачок).

\subsection{Удобная форма уравнения на число частиц}

Запишем уравнение~\eqref{eq:N_bose} в удобном безразмерном виде. Введём квантовый объём $V_Q = \left(\frac{2\pi\hbar^2}{mT}\right)^{3/2}$ --- объём фазовой ячейки порядка $\hbar^3$, делённый на импульс (тепловой). Тогда $A = \frac{g\,V}{2\pi^2}\left(\frac{2m}{\hbar^2}\right)^{3/2}\!\cdot\frac{1}{2}$, и уравнение на $N$ принимает вид:
\begin{equation}
    n V_Q = \frac{2}{\sqrt{\pi}}\,I_{1/2}(y) \equiv f(y)
    \label{eq:nVQ}
\end{equation}
где $n = N/V$ --- концентрация. Функция $f(y)$ монотонно убывает от:
\begin{equation}
    f(0) = \frac{2}{\sqrt{\pi}}\,I_{1/2}(0) = \frac{2}{\sqrt{\pi}}\cdot\frac{\sqrt{\pi}}{2}\cdot\zeta\!\left(\tfrac{3}{2}\right) = \zeta\!\left(\tfrac{3}{2}\right) \approx 2{,}612
\end{equation}
до нуля при $y \to \infty$. Уравнение~\eqref{eq:nVQ} при фиксированной концентрации $n$ и убывающем с ростом $T$ параметре $V_Q \propto T^{-3/2}$ определяет $y$ как функцию температуры.

% ===== КОНЕЦ СЕКЦИИ 1 =====


% ===== СЕКЦИЯ 2: Бозе-Эйнштейновская конденсация =====
\section{Бозе-Эйнштейновская конденсация}

\subsection{Когда уравнение перестаёт иметь решение}

Ключевое наблюдение: функция $f(y)$ ограничена сверху значением $f(0) = \zeta(3/2) \approx 2{,}612$.

При высоких температурах $V_Q$ мало, $nV_Q \ll 1$, и уравнение~\eqref{eq:nVQ} имеет решение при некотором $y \gg 1$ (классический режим). При понижении температуры $V_Q$ растёт, $nV_Q$ увеличивается, и $y$ уменьшается. При некоторой температуре $T_c$ получаем $nV_Q(T_c) = \zeta(3/2)$, то есть $y = 0$ и $\mu = 0$.

При $T < T_c$ величина $nV_Q > \zeta(3/2)$, и уравнение~\eqref{eq:nVQ} не имеет решения при $y > 0$. Это означает, что наша исходная формула для $N$, содержащая только непрерывный спектр (интеграл по $d\epsilon$), перестаёт быть применимой.

Физически это означает: часть частиц «не помещается» в возбуждённые состояния непрерывного спектра. Они вынуждены конденсироваться на нулевом уровне энергии $\epsilon = 0$ --- это и есть \textbf{бозе-эйнштейновская конденсация}.

\subsection{Критическая температура}

Критическая температура $T_c$ определяется из условия $y = 0$, $\mu = 0$:
\begin{equation}
    n\,V_Q(T_c) = \zeta\!\left(\tfrac{3}{2}\right) \approx 2{,}612
\end{equation}
Подставляя $V_Q = \left(\frac{2\pi\hbar^2}{mT_c}\right)^{3/2}$:
\begin{equation}
    \boxed{T_c = \frac{2\pi\hbar^2}{m}\left(\frac{n}{\zeta(3/2)}\right)^{\!2/3}}
    \label{eq:Tc}
\end{equation}

Заметим, что $T_c \propto n^{2/3}$: чем плотнее газ, тем выше температура конденсации. Это интуитивно понятно: конденсация происходит, когда квантовые волновые пакеты частиц начинают перекрываться.

\subsection{Оценка для гелия-4}

Сделаем численную оценку для жидкого гелия-4. Масса атома $m = 4\,m_u \approx 6{,}64\cdot10^{-27}$ кг, молярный объём $V_\text{мол} \approx 27{,}6$ см³/моль, откуда концентрация:
\begin{equation}
    n = \frac{N_A}{V_\text{мол}} \approx \frac{6{,}02\cdot10^{23}}{27{,}6\cdot10^{-6}\text{ м}^3} \approx 2{,}18\cdot10^{28}\text{ м}^{-3}
\end{equation}

Подставляя в формулу~\eqref{eq:Tc}: $T_c^{(\text{расч})} \approx 3{,}13$ К.

Экспериментальное значение: $\lambda$-переход в гелии-4 наблюдается при $T_\lambda \approx 2{,}19$ К. Совпадение неплохое --- учитывая, что жидкий гелий совсем не является идеальным газом! Различие объясняется межатомными взаимодействиями, которые мы не учли.

\subsection{Картина конденсации: при $T < T_c$}

\begin{figure}[h!]
    \centering
    %[Место для графика: Химический потенциал $\mu$ как функция температуры $T$. При $T > T_c$: $\mu < 0$, монотонно убывает вниз. При $T = T_c$: $\mu = 0$. При $T < T_c$: горизонтальная прямая $\mu = 0$. Таймкод: 18:00--20:00]
    \caption{Химический потенциал идеального бозе-газа как функция температуры при фиксированной концентрации.}
\end{figure}

При $T < T_c$ химический потенциал остаётся равным нулю: $\mu = 0$ при всех $T \leq T_c$. Это вынужденное условие: если $\mu = 0$, то все частицы, не умещающиеся в непрерывный спектр, занимают основной уровень $\epsilon = 0$.

Разделим общее число частиц на две части: $N = N_0 + N'$, где $N_0$ --- число частиц в конденсате (на уровне $\epsilon = 0$), $N'$ --- число частиц в возбуждённых состояниях.

При $\mu = 0$ число частиц в возбуждённых состояниях равно:
\begin{equation}
    N' = A\,T^{3/2}\,I_{1/2}(0) = A\,T^{3/2}\cdot\frac{\sqrt{\pi}}{2}\,\zeta\!\left(\tfrac{3}{2}\right)
\end{equation}

При $T = T_c$ это равно $N$:
\begin{equation}
    N = A\,T_c^{3/2}\cdot\frac{\sqrt{\pi}}{2}\,\zeta\!\left(\tfrac{3}{2}\right)
\end{equation}

Разделив первое на второе:
\begin{equation}
    \frac{N'}{N} = \left(\frac{T}{T_c}\right)^{3/2}
\end{equation}

Следовательно, число частиц в конденсате:
\begin{equation}
    \boxed{N_0 = N\left[1 - \left(\frac{T}{T_c}\right)^{3/2}\right]}
    \label{eq:N0}
\end{equation}

При $T \to 0$ все частицы конденсируются ($N_0 \to N$). При $T \to T_c$ конденсат исчезает ($N_0 \to 0$).

\subsection{Исторические замечания}

Бозе-эйнштейновская конденсация --- одно из самых замечательных квантовых явлений. Исторически ситуация была следующей. Предсказали её Бозе и Эйнштейн в 1924--1925 годах. Долгое время единственным кандидатом на роль бозе-конденсата оставался жидкий гелий-4 ниже $\lambda$-точки. Однако жидкий гелий --- сильно взаимодействующая система, и чистую конденсацию идеального газа наблюдать не удавалось.

Лишь в 1995 году удалось получить бозе-эйнштейновский конденсат в парах щелочных металлов (рубидий, натрий) --- при температурах порядка $10^{-7}$ К и концентрациях, на много порядков ниже, чем в обычных газах. За это открытие Корнелл, Виман и Кеттерле получили Нобелевскую премию по физике в 2001 году.

% ===== КОНЕЦ СЕКЦИИ 2 =====


% ===== СЕКЦИЯ 3: Энергия бозе-газа =====
\section{Энергия идеального бозе-газа}

\subsection{Общая формула}

Средняя энергия бозе-газа записывается аналогично числу частиц:
\begin{equation}
    E = \int_0^\infty \frac{\epsilon\,\nu(\epsilon)}{e^{(\epsilon-\mu)/T}-1}\, d\epsilon = A\int_0^\infty \frac{\epsilon^{3/2}}{e^{(\epsilon-\mu)/T}-1}\, d\epsilon
\end{equation}
В безразмерных переменных ($x = \epsilon/T$, $y = -\mu/T$):
\begin{equation}
    E = A\,T^{5/2}\, I_{3/2}(y)
    \label{eq:E_bose}
\end{equation}

\subsection{Высокотемпературный предел ($T \gg T_c$, $y \gg 1$)}

При больших $y$ используем $I_{3/2}(y) \approx \Gamma(5/2)\,e^{-y} = \frac{3\sqrt{\pi}}{4}\,e^{-y}$. Аналогично, $I_{1/2}(y) \approx \frac{\sqrt{\pi}}{2}\,e^{-y}$. Поэтому:
\begin{equation}
    \frac{E}{N} = \frac{A\,T^{5/2}\,I_{3/2}(y)}{A\,T^{3/2}\,I_{1/2}(y)} = T\cdot\frac{I_{3/2}(y)}{I_{1/2}(y)} \approx T\cdot\frac{(3\sqrt{\pi}/4)e^{-y}}{(\sqrt{\pi}/2)e^{-y}} = \frac{3}{2}\,T
\end{equation}

Итак, при $T \gg T_c$: $E = \frac{3}{2}NT$ --- классический результат.

\subsection{Низкотемпературный режим ($T < T_c$, $y = 0$)}

При $T < T_c$ химический потенциал обращается в нуль: $\mu = 0$, то есть $y = 0$. Только возбуждённые частицы вносят вклад в энергию (частицы в конденсате имеют нулевую энергию). Подставляя $y = 0$ в формулу~\eqref{eq:E_bose}:
\begin{equation}
    E = A\,T^{5/2}\,I_{3/2}(0) = A\,T^{5/2}\cdot\frac{3\sqrt{\pi}}{4}\,\zeta\!\left(\tfrac{5}{2}\right)
\end{equation}

Выразим это через $N$ и $T_c$. Из определения $T_c$:
\begin{equation}
    N = A\,T_c^{3/2}\cdot\frac{\sqrt{\pi}}{2}\,\zeta\!\left(\tfrac{3}{2}\right)
\end{equation}
Поэтому:
\begin{equation}
    E = N\,T\left(\frac{T}{T_c}\right)^{3/2}\cdot\frac{I_{3/2}(0)}{I_{1/2}(0)} = N\,T\left(\frac{T}{T_c}\right)^{3/2}\cdot\frac{(3\sqrt{\pi}/4)\,\zeta(5/2)}{(\sqrt{\pi}/2)\,\zeta(3/2)}
\end{equation}

Вычислим коэффициент:
\begin{equation}
    \frac{I_{3/2}(0)}{I_{1/2}(0)} = \frac{3}{2}\cdot\frac{\zeta(5/2)}{\zeta(3/2)} \approx \frac{3}{2}\cdot\frac{1{,}341}{2{,}612} \approx 0{,}770
\end{equation}

Итого:
\begin{equation}
    \boxed{E = 0{,}770\;N\,T\left(\frac{T}{T_c}\right)^{3/2}, \quad T < T_c}
    \label{eq:E_bose_condensate}
\end{equation}

% ===== КОНЕЦ СЕКЦИИ 3 =====


% ===== СЕКЦИЯ 4: Теплоёмкость бозе-газа =====
\section{Теплоёмкость идеального бозе-газа}

\subsection{Теплоёмкость при $T < T_c$}

Дифференцируем формулу~\eqref{eq:E_bose_condensate} по температуре (при $T < T_c$ формула явная, $T_c$ фиксировано):
\begin{equation}
    C_V = \frac{dE}{dT} = 0{,}770\;N\,T_c^{-3/2}\cdot\frac{5}{2}\,T^{3/2} = \frac{5}{2}\cdot 0{,}770\;N\left(\frac{T}{T_c}\right)^{3/2} \approx 1{,}925\;N\left(\frac{T}{T_c}\right)^{3/2}
\end{equation}

Таким образом, при $T < T_c$:
\begin{equation}
    C_V \approx 1{,}925\;N\left(\frac{T}{T_c}\right)^{3/2}
\end{equation}

\subsection{Значение теплоёмкости в точке $T_c$}

В точке конденсации:
\begin{equation}
    C_V(T_c^-) \approx 1{,}925\,N
\end{equation}

Это \emph{больше} классического значения $\frac{3}{2}N$! Физически это объяснимо: около $T_c$ бозоны накапливаются на низких уровнях, и система становится особенно чувствительной к изменению температуры.

\subsection{Поведение теплоёмкости при $T > T_c$}

При $T > T_c$ теплоёмкость убывает обратно к классическому значению $\frac{3}{2}N$. Точный расчёт, учитывающий зависимость $\mu(T)$ при $T > T_c$, показывает, что $C_V$ при прохождении через $T_c$ не имеет разрыва, однако её производная по температуре ($dC_V/dT$) испытывает скачок --- \emph{излом}.

\begin{figure}[h!]
    \centering
    %[Место для графика: Теплоёмкость $C_V$ как функция $T/T_c$. При $T < T_c$: кривая пропорциональна $(T/T_c)^{3/2}$, монотонно растёт. В точке $T_c$: максимум $C_V \approx 1{,}925\,N$ с изломом (производная скачком меняется). При $T > T_c$: монотонно убывает к $\frac{3}{2}N$. Таймкод: 32:00--35:00]
    \caption{Теплоёмкость идеального бозе-газа в зависимости от температуры. Излом в точке $T_c$ --- признак фазового перехода.}
\end{figure}

Этот излом --- признак фазового перехода. В отличие от ферми-газа, где теплоёмкость плавно линейна по $T$, в бозе-газе при $T_c$ происходит качественное изменение: часть частиц образует конденсат и «выпадает» из термодинамики возбуждённых состояний.

% ===== КОНЕЦ СЕКЦИИ 4 =====


% ===== СЕКЦИЯ 5: Уравнение состояния и изотермы бозе-газа =====
\section{Уравнение состояния и изотермы идеального бозе-газа}

\subsection{Связь давления с энергией}

Как мы установили ещё при изучении ферми-газа, для любого идеального невзаимодействующего нерелятивистского газа справедливо:
\begin{equation}
    PV = \frac{2}{3}\,E
\end{equation}
Это соотношение вытекает из квадратичного закона дисперсии $\epsilon = p^2/(2m)$ и не зависит от статистики. Поэтому оно справедливо и для бозе-газа.

\subsection{Уравнение состояния при $T > T_c$}

При высоких температурах бозе-газ переходит в классический: $E = \frac{3}{2}NT$, откуда:
\begin{equation}
    P = \frac{2E}{3V} = \frac{NT}{V} = nT
\end{equation}
Получаем уравнение состояния идеального классического газа.

\subsection{Уравнение состояния при $T < T_c$}

При $T < T_c$ подставляем формулу~\eqref{eq:E_bose_condensate}:
\begin{equation}
    P = \frac{2E}{3V} = \frac{2\cdot 0{,}770}{3}\,\frac{N}{V}\,T\left(\frac{T}{T_c}\right)^{3/2}
\end{equation}
Поскольку $T_c \propto n^{2/3}$, давление $P$ при $T < T_c$ пропорционально $T^{5/2}$ и \emph{не зависит от объёма}:
\begin{equation}
    P \propto T^{5/2}, \quad T < T_c
\end{equation}

Более точно: можно показать, что:
\begin{equation}
    P = \frac{2}{3}\,A\,T^{5/2}\,I_{3/2}(0) \approx 0{,}0851\,\frac{m^{3/2}}{\hbar^3}\,T^{5/2}
\end{equation}
(это число получается подстановкой числовых значений констант). Давление зависит только от температуры, но не от объёма.

\subsection{Нарушение условия механической устойчивости}

Условие механической устойчивости термодинамической системы: $\left(\partial P/\partial V\right)_T < 0$. В нашем случае при $T < T_c$:
\begin{equation}
    \left(\frac{\partial P}{\partial V}\right)_T = 0
\end{equation}

Производная равна нулю, что формально нарушает условие устойчивости. Это является следствием идеализации: в идеальном газе нет взаимодействия между частицами. Любое, сколь угодно слабое отталкивание между бозонами привело бы к конечному наклону изотерм и восстановило бы механическую устойчивость.

\subsection{Изотермы бозе-газа на диаграмме $P$--$V$}

\begin{figure}[h!]
    \centering
    %[Место для графика: Изотермы на диаграмме $P$--$V$. При $T > T_c$: гиперболические кривые $P \propto 1/V$ (классика с квантовыми поправками). Граничная кривая (локус точек перехода в конденсат): $P \propto V^{-5/3}$, то есть $PV^{5/3} = \text{const}$. При $T < T_c$ (область внутри граничной кривой): горизонтальная линия $P = P(T) = \text{const}$ при больших $V$ (конденсат). При малых $V$ (за границей конденсации): классическая кривая. Таймкод: 40:00--44:00]
    \caption{Изотермы идеального бозе-газа. Горизонтальный участок соответствует области бозе-конденсации.}
\end{figure}

Описание изотерм:
\begin{itemize}
    \item При $T > T_c$: при фиксированной температуре $P \propto 1/V$ (классика). При уменьшении $V$ (увеличении $n$) при некотором $V = V_c(T)$ система входит в режим конденсации.
    \item При $V > V_c(T)$: горизонтальная линия $P = P(T)$, давление не меняется при изменении объёма.
    \item Граничная кривая (геометрическое место точек начала конденсации) на диаграмме $P$--$V$: поскольку $P \propto T^{5/2}$ и $V_c \propto T^{-3/2}$, исключая $T$, получаем:
    \begin{equation}
        P\cdot V^{5/3} = \text{const}
    \end{equation}
    Эта кривая напоминает изотерму, но с показателем $5/3$, а не $1$. Поведение аналогично жидко-газовому фазовому переходу (ср. с изотермой Ван-дер-Ваальса).
\end{itemize}

% ===== КОНЕЦ СЕКЦИИ 5 =====


% ===== СЕКЦИЯ 6: Первый закон для открытой системы и термодинамические потенциалы =====
\section{Термодинамические потенциалы для открытой системы}

Теперь переходим к новой большой теме --- термодинамические потенциалы. Всё, что мы делали раньше, вело нас к этому.

\subsection{Первый закон для открытой системы}

До сих пор мы рассматривали закрытые системы с фиксированным числом частиц $N$. Теперь обобщим на открытые системы, в которых $N$ может меняться.

В этом случае первый закон термодинамики принимает вид:
\begin{equation}
    \boxed{dE = T\,dS - P\,dV + \mu\,dN}
    \label{eq:first_law_open}
\end{equation}

Здесь $\mu$ --- химический потенциал, который мы хорошо знаем из статистической физики. Термодинамически $\mu$ --- это работа по добавлению одной частицы в систему при фиксированных $S$ и $V$.

Формула~\eqref{eq:first_law_open} говорит нам, что внутренняя энергия $E$ является функцией трёх переменных: $E = E(S, V, N)$. Все три переменные --- экстенсивные (масштабируются с размером системы).

\subsection{Термодинамические потенциалы: идея Лежандрова преобразования}

На практике бывает удобнее работать не с переменными $(S, V, N)$, а с другими наборами переменных. Например, в лаборатории часто удобнее контролировать температуру $T$, а не энтропию $S$; давление $P$, а не объём $V$. Переход от одного набора переменных к другому осуществляется с помощью \textbf{преобразования Лежандра}.

Идея: если $E = E(S, V, N)$, то из $dE = TdS - PdV + \mu\,dN$ видно, что:
\begin{equation}
    T = \left(\frac{\partial E}{\partial S}\right)_{V,N}, \quad -P = \left(\frac{\partial E}{\partial V}\right)_{S,N}, \quad \mu = \left(\frac{\partial E}{\partial N}\right)_{S,V}
\end{equation}

Чтобы перейти к функции переменных $(T, V, N)$, нужно заменить $S \to T$: для этого вычитаем из $E$ произведение $T\cdot S$.

\subsection{Свободная энергия Гельмгольца}

Определим свободную энергию Гельмгольца:
\begin{equation}
    F = E - T\,S
\end{equation}
Вычислим полный дифференциал:
\begin{equation}
    dF = dE - T\,dS - S\,dT = (T\,dS - P\,dV + \mu\,dN) - T\,dS - S\,dT
\end{equation}
\begin{equation}
    \boxed{dF = -S\,dT - P\,dV + \mu\,dN}
    \label{eq:dF}
\end{equation}

Итак, $F = F(T, V, N)$ --- функция температуры, объёма и числа частиц. Это наиболее часто используемый потенциал в статистической физике, потому что именно $T$, $V$, $N$ легко контролировать.

Из формулы~\eqref{eq:dF}:
\begin{equation}
    S = -\left(\frac{\partial F}{\partial T}\right)_{V,N}, \quad P = -\left(\frac{\partial F}{\partial V}\right)_{T,N}, \quad \mu = \left(\frac{\partial F}{\partial N}\right)_{T,V}
\end{equation}

\subsection{Энтальпия}

Другое преобразование: перейдём от объёма $V$ к давлению $P$. Для этого прибавим к $E$ произведение $P\cdot V$:
\begin{equation}
    H = E + P\,V
\end{equation}
\begin{equation}
    dH = dE + P\,dV + V\,dP = T\,dS + V\,dP + \mu\,dN
\end{equation}
\begin{equation}
    \boxed{dH = T\,dS + V\,dP + \mu\,dN}
    \label{eq:dH}
\end{equation}

Энтальпия $H = H(S, P, N)$ удобна при постоянном давлении (например, химические реакции при атмосферном давлении).

\subsection{Потенциал Гиббса}

Третье преобразование: заменим одновременно $S \to T$ и $V \to P$. Получим потенциал Гиббса (или функцию Гиббса):
\begin{equation}
    \Phi = E - T\,S + P\,V = F + PV = H - TS
\end{equation}
\begin{equation}
    d\Phi = dE - T\,dS - S\,dT + P\,dV + V\,dP = -S\,dT + V\,dP + \mu\,dN
\end{equation}
\begin{equation}
    \boxed{d\Phi = -S\,dT + V\,dP + \mu\,dN}
    \label{eq:dPhi}
\end{equation}

Потенциал Гиббса $\Phi = \Phi(T, P, N)$ удобен при постоянных $T$ и $P$ --- именно в таких условиях обычно происходят фазовые переходы.

\subsection{Теорема Эйлера: потенциал Гиббса и химический потенциал}

Потенциал Гиббса обладает замечательным свойством. Все переменные в $\Phi(T, P, N)$ таковы, что только $N$ является экстенсивной переменной ($T$ и $P$ --- интенсивные). Следовательно, по теореме Эйлера об однородных функциях:
\begin{equation}
    \Phi(T, P, \lambda N) = \lambda\,\Phi(T, P, N)
\end{equation}
Дифференцируя по $\lambda$ и полагая $\lambda = 1$:
\begin{equation}
    \boxed{\Phi = \mu(T, P)\cdot N}
    \label{eq:Phi_muN}
\end{equation}

Это фундаментальный результат: потенциал Гиббса всей системы из $N$ частиц равен $\mu N$, где $\mu = \mu(T, P)$ --- химический потенциал, зависящий только от интенсивных переменных.

\subsection{Теорема малых добавок: химический потенциал как частная производная}

Из всего написанного видно, что химический потенциал допускает четыре эквивалентных представления:
\begin{equation}
    \boxed{\mu = \left(\frac{\partial E}{\partial N}\right)_{S,V} = \left(\frac{\partial H}{\partial N}\right)_{S,P} = \left(\frac{\partial F}{\partial N}\right)_{T,V} = \left(\frac{\partial \Phi}{\partial N}\right)_{T,P}}
    \label{eq:mu_four}
\end{equation}

Каждое из этих выражений --- это работа по добавлению одной частицы при соответствующих фиксированных параметрах. Это и называется \textbf{теоремой малых добавок} (или теоремой об инвариантности химического потенциала).

% ===== КОНЕЦ СЕКЦИИ 6 =====


% ===== СЕКЦИЯ 7: Условия минимальности термодинамических потенциалов =====
\section{Условия минимальности термодинамических потенциалов}

\subsection{Второй закон термодинамики}

Перейдём к вопросу о равновесии. Второй закон термодинамики в форме, пригодной для рассуждений:

Для равновесного процесса:
\begin{equation}
    dS = \frac{\delta Q}{T}
\end{equation}
Для любого реального (в том числе необратимого) процесса:
\begin{equation}
    dS \geq \frac{\delta Q}{T_\text{env}}
\end{equation}
где $T_\text{env}$ --- температура нагревателя, термостата (внешнего источника тепла). Знак равенства достигается лишь в квазистатическом обратимом процессе.

\subsection{Состояние частичного равновесия}

Рассмотрим следующую ситуацию. Имеется изолированная система, разделённая перегородкой на две части. Каждая часть находится в состоянии локального равновесия со своей температурой: левая подсистема при температуре $T_1$, правая --- при $T_2$, причём $T_1 \neq T_2$. Это называется \textbf{состоянием частичного равновесия}: локально всё хорошо, но глобально система не в равновесии.

Теперь уберём перегородку (или сделаем её теплопроводящей). Произойдёт спонтанный необратимый процесс выравнивания температур: тепло потечёт от горячей подсистемы к холодной, пока не установится общая температура $T_1 = T_2 = T$.

Вопрос: как ведут себя термодинамические потенциалы в ходе этого процесса? Что минимизируется, а что максимизируется?

\subsection{Максимизация энтропии в изолированной системе}

Для изолированной системы ($\delta Q = 0$ снаружи, $V = \text{const}$, $N = \text{const}$) второй закон говорит:
\begin{equation}
    dS \geq 0
\end{equation}
Энтропия в изолированной системе никогда не убывает. В состоянии равновесия она достигает максимума. Неравновесный спонтанный переход от $T_1 \neq T_2$ к $T_1 = T_2$ сопровождается ростом суммарной энтропии.

\subsection{Что минимизируется при других условиях}

Нас интересует, что происходит при других фиксированных параметрах. Это будет детально разобрано на следующей лекции. Забегу вперёд, чтобы у вас была общая картина.

Оказывается:
\begin{itemize}
    \item При фиксированных $T$ и $V$: в ходе спонтанного процесса к равновесию \textbf{минимизируется свободная энергия Гельмгольца} $F$.
    \item При фиксированных $T$ и $P$: в ходе спонтанного процесса к равновесию \textbf{минимизируется потенциал Гиббса} $\Phi$.
\end{itemize}

Это делает $F$ и $\Phi$ центральными объектами при описании фазовых переходов: фаза, соответствующая меньшему $F$ или $\Phi$ при данных условиях, является термодинамически устойчивой.

На сегодня наша лекция заканчивается. В следующий раз мы строго докажем условия минимальности потенциалов и применим их к задачам о фазовых переходах.

% ===== КОНЕЦ СЕКЦИИ 7 =====


\end{document}
